\documentclass[10pt]{article}
\usepackage{formal}

% Packages for ToC and formatting
\usepackage{amsmath}
\usepackage{amssymb}
\usepackage{mathrsfs}
\usepackage{geometry}
\usepackage{unicode-math}
\usepackage{hyperref}
\usepackage{cleveref}
\usepackage{tocbibind}

% Document setup
\title{WebSocket Machine Formal Specification}
\author{Zhifeng Zhang}

\begin{document}

\maketitle
\tableofcontents
\newpage

%%--------------------------------------------------------------------------------
%% Core module
%%--------------------------------------------------------------------------------
\section{Core Machine Model}

\subsection{Primary Definition}
A WebSocket Machine is formally defined as a 7-tuple:
\[
\mathfrak{M} = (S, E, \delta, s_0, C, \gamma, F)
\]
Where
\begin{itemize}[leftmargin=2em]
\item $S$ is the finite set of states
\item $E$ is the set of events
\item $\delta: S \times E \rightarrow S \times \gamma$ is the transition function
\item $s_0 \in S$ is the initial state
\item $C$ is the context triple
\item $\gamma$ is the set of actions 
\item $F \subseteq S$ is the set of final states
\end{itemize}

%%-------- States ($S$) ---------------------------------------------------------
\subsection{States}
The state set is defined as:
\[
S = \{s_i \mid i=1,2,...,n;\ n=6\}
\]
where
\begin{eqnarray*}
s_1 &=& \text{Disconnected} \\
s_2 &=& \text{Connecting} \\
s_3 &=& \text{Connected} \\
s_4 &=& \text{Reconnecting} \\
s_5 &=& \text{Disconnecting} \\
s_6 &=& \text{Terminated}
\end{eqnarray*}

Furthermore, the initial state $s_0$ is defined to be $s_1$ and the final state
set $F=\{s_6\}$ or $F=\{\text{Terminated}\}$.

%%-------- Events ($E$) ---------------------------------------------------------
\subsection{Events}
The event set is defined as:
\[
E = \{e_i \mid i=1,2,...,m;\ m=12\}
\]
where
\begin{eqnarray*}
e_1 &=& \text{CONNECT} \\
e_2 &=& \text{DISCONNECT} \\
e_3 &=& \text{OPEN} \\
e_4 &=& \text{CLOSE} \\
e_5 &=& \text{ERROR} \\
e_6 &=& \text{MESSAGE} \\
e_7 &=& \text{PING} \\
e_8 &=& \text{PONG} \\
e_9 &=& \text{TIMEOUT} \\
e_{10} &=& \text{RETRY} \\
e_{11} &=& \text{ABORT} \\
e_{12} &=& \text{TERMINATE}
\end{eqnarray*}

%%-------- Context ($C$) ---------------------------------------------------------
\subsection{Context}
The context is defined as a triple:
$$C = (P, V, T)$$
where
\begin{itemize}[leftmargin=2em]
\item $P$ is the set of primary connection properties
  $$P = \{\text{url}, \text{protocols}, \text{socket}, \text{status}, \text{readyState}\}$$
\item $V$ is the set of metric values
  $$V = \{\text{messagesSent}, \text{messagesReceived}, \text{reconnectAttempts},
  \text{bytesSent}, \text{bytesReceived}\}$$
\item $T$ is the set of timing properties
  $$T = \{\text{connectTime}, \text{disconnectTime}, \text{lastPingTime},
  \text{lastPongTime}, \text{windowStart}\}$$
\end{itemize}

%%-------- Actions ($\gamma$) -----------------------------------------------------
\subsection{Actions}
\addtocontents{toc}{\protect\setcounter{tocdepth}{2}}

Actions are pure functions that transform context:
$$\gamma: C \times E \rightarrow C$$

Core action set:
$$\Gamma = \{\gamma_i \mid i=1,2,...,p;\ p=11\}$$

Key actions include:

\begin{eqnarray*}
\gamma_1(c, e_{\text{\tiny CONNECT}}) &=& c' \text{ where } c'.\text{url} = e.\text{url} \\
\gamma_2(c) &=& c' \text{ where } c'.\text{reconnectAttempts} = 0 \\
\gamma_3(c, e_{\text{\tiny ERROR}}) &=& c' \text{ where } c'.\text{error} = e.\text{error} \\
\gamma_4(c, e_{\text{\tiny MESSAGE}}) &=& c' \text{ where } c'.\text{messagesReceived} += 1 \\
\gamma_5(c, e_{\text{\tiny SEND}}) &=& c' \text{ where } c'.\text{messagesSent} += 1
\end{eqnarray*}

%%-------- Transition Function ($\delta$) -----------------------------------------
\subsection{Transition Function}

The transition function maps states and events to new states and actions:
$$\delta: S \times E \rightarrow S \times \mathfrak{P}(\Gamma)$$

Key transitions include:
\begin{eqnarray*}
\delta(s_{\text{\tiny Disconnected}}, e_{\text{\tiny CONNECT}}) &=& (s_{\text{\tiny Connecting}}, \{\gamma_1, \gamma_{10}\}) \\
\delta(s_{\text{\tiny Connecting}}, e_{\text{\tiny OPEN}}) &=& (s_{\text{\tiny Connected}}, \{\gamma_2\}) \\
\delta(s_{\text{\tiny Connected}}, e_{\text{\tiny ERROR}}) &=& (s_{\text{\tiny Reconnecting}}, \{\gamma_3, \gamma_9\})
\end{eqnarray*}


%%--------------------------------------------------------------------------------
%% Support module
%%--------------------------------------------------------------------------------
\section{Support Systems}

%%-------- Type System ($\mathfrak{T}$) -------------------------------------------
\subsection{Type System ($\mathfrak{T}$)}

A Type System is defined as:
$$\mathfrak{T} = (B, C, V)$$

where:
- $B$: Base types $\{\text{String}, \text{Number}, \text{Boolean}, \text{Null}, \text{Undefined}\}$
- $C$: Composite types (Record, Array, Union, Intersection)
- $V$: Validation functions

%%-------- Guards System ($\mathfrak{G}$) -----------------------------------------
\subsection{Guards System ($\mathfrak{G}$)}

A Guards System is defined as:
\[
\mathfrak{G} = (P, \Gamma, \Psi, \Omega)
\]

where:
\begin{itemize}[leftmargin=2em]
\item $P$: Predicate functions
\item $\Gamma$: Guard compositions
\item $\Psi$: Guard evaluation context
\item $\Omega$: Guard operators
\end{itemize}

Guard composition laws:
\begin{eqnarray}
\gamma_{\text{\tiny AND}}(p_1, p_2)(c, e) &=& p_1(c, e) \wedge p_2(c, e) \\
\gamma_{\text{\tiny OR}}(p_1, p_2)(c, e) &=& p_1(c, e) \vee p_2(c, e) \\
\gamma_{\text{\tiny NOT}}(p)(c, e) &=& \neg p(c, e)
\end{eqnarray}

%%-------- Error System ($\varepsilon$) -----------------------------------------
\subsection{Error System}

An Error System is defined as:
$$\varepsilon = (K, R, H)$$
where
\begin{itemize}[leftmargin=2em]
\item $K$: Error categories
\item $R$: Recovery strategies
\item $H$: Error history
\end{itemize}

Recovery strategy function:
$$R(k, n) = (m, b, t)$$
where
\begin{itemize}[leftmargin=2em]
\item $m$: Maximum retry attempts
\item $b$: Backoff factor
\item $t$: Timeout duration
\end{itemize}

%%-------- Resource Management ($\mathfrak{R}$) -----------------------------------
\subsection{Resource Management ($\mathfrak{R}$)}

A Resource Management System $\mathfrak{R}$ is defined as:
$$\mathfrak{R} = (L, A, M)$$
where
\begin{itemize}[leftmargin=2em]
\item $L$: Lifecycle states $\{\text{free}, \text{acquired}, \text{failed}\}$
\item $A$: Acquisition operations
\item $M$: Monitoring functions
\end{itemize}

%%-------- Rate Limiting ($\rho$) -----------------------------------------------
\subsection{Rate Limiting}

A Rate Limiting System $\rho$ is defined as:
$$\rho = (W, Q, \lambda)$$
where
\begin{itemize}[leftmargin=2em]
\item $W$: Time window function
\item $Q$: Message queue
\item $\lambda$: Rate limiting function
\end{itemize}

Window function:
$$W(t) = [t - w, t]$$

Rate limiting function:
$$
\lambda(m, W) = \begin{cases}
\text{accept} & \text{if } |M(W)| < n \\
\text{queue} & \text{if } |Q| < q_{\text{\tiny max}} \\
\text{reject} & \text{otherwise}
\end{cases}
$$

%%-------- Health Monitoring ($\mathfrak{H}$) -------------------------------------
\subsection{Health Monitoring ($\mathfrak{H}$)}

A Health Monitoring System is defined as:
$$\mathfrak{H} = (\Pi, \Delta, \Phi)$$
where
\begin{itemize}[leftmargin=2em]
\item $\Pi$: Health probe system
\item $\Delta$: Health metrics $(L, R, U)$
\item $\Phi$: Health state function
\end{itemize}

Health state function:
$$
\Phi(\Delta) = \begin{cases}
\text{healthy} & \text{if } R \geq r_{\text{\tiny threshold}} \wedge \bar{L} \leq l_{\text{\tiny threshold}} \\
\text{degraded} & \text{if } R \geq r_{\text{\tiny critical}} \wedge \bar{L} \leq l_{\text{\tiny critical}} \\
\text{unhealthy} & \text{otherwise}
\end{cases}
$$

%%-------- Metrics Collection ($\mathfrak{M}$) -----------------------------------
\subsection{Metrics Collection ($\mathfrak{M}$)}

A Metrics Collection System is defined as:
$$\mathfrak{M} = (D, \Sigma, \Theta, \Phi)$$
where
\begin{itemize}[leftmargin=2em]
\item $D$: Metric data points
\item $\Sigma$: Collection strategies
\item $\Theta$: Temporal aggregation
\item $\Phi$: Processing functions
\end{itemize}

%%-------- Testing Verification ($\mathfrak{V}$) ----------------------------------
\subsection{Testing Verification ($\mathfrak{V}$)}

A Testing Verification System is defined as:
$$\mathfrak{V} = (P, I, T)$$
where
\begin{itemize}[leftmargin=2em]
\item $P$: Property verification functions
\item $I$: Invariant checks
\item $T$: Transition tests
\end{itemize}

%%--------------------------------------------------------------------------------
%% Support module
%%--------------------------------------------------------------------------------
\section{Integration Properties}

%%-------- System Composition Function -----------------------------------------
\subsection{System Composition Function}

The complete system is composed as:
$$\mathfrak{S} = (\mathfrak{T}, \varepsilon, \mathfrak{R}, \rho, \mathfrak{H}, \mathfrak{V})$$

%%-------- System Invariants ---------------------------------------------------
\subsection{System Invariants}

\begin{enumerate}[leftmargin=2em]
\item Type Safety:
   $$\forall x, \forall a,b \in \mathfrak{S}: V_a(x) \wedge V_b(x)$$
\item Error Commutativity:
   $$\varepsilon_a \circ \varepsilon_b = \varepsilon_b \circ \varepsilon_a$$
\item Resource Exclusivity:
   $$\forall a,b \in \mathfrak{S}: \mathfrak{R}_a \cap \mathfrak{R}_b = \emptyset$$
\item Rate Limit Monotonicity:
   $$t_1 \leq t_2 \implies W(t_1) \subseteq W(t_2)$$
\item Health Continuity:
   $$\forall t_1,t_2: |t_1-t_2| < \epsilon \implies |\Phi(\Delta(t_1)) - \Phi(\Delta(t_2))| < \delta$$
\end{enumerate}

%%-------- System Properties ---------------------------------------------------
\subsection{Global Properties}

\begin{enumerate}[leftmargin=2em]
\item Determinism:
   $$\forall s \in S, e \in E: |\delta(s,e)| \leq 1$$
\item Reachability:
   $$\forall s \in S: \exists \text{ path } p: s_0 \xrightarrow{*} s$$
\item Safety:
   $$\forall s \in S - F: \exists e \in E: \delta(s,e) \neq \emptyset$$
\item Liveness:
   $$\forall s \in S: \exists \text{ path } p: s \xrightarrow{*} f \text{ where } f \in F$$
\end{enumerate}
\end{document}